Emmy haar pubertijd
gaf haar onzekerheid

over de succeskans
van haar temmende sjans.

Toen werd ze vaak gepest.
Haar eer steeds uitgetest.

Wat ze haar aandeden,
kon haar niet downgraden

tot hun verlaagd niveau.
Ze liet het dus maar zo.

Omdat ze steeds aannam
dat het vanzelf goed kwam.

Haar tijd in het klooster
verbrak het pestrooster,

tot haar kerkgenoten
wat anders besloten.

Haar werd, zonder weerwoord,
zonder te zijn gehoord,

een misstap verweten
en eruit gesmeten.

Zo werd ze op een dag
bang als ze pesters zag.

Op zaterdagavond
hingen ze meestal rond

Emmy’s ouderlijk huis.
Ze bleef dan het liefst thuis.

Wat niet een optie was.
Ze sloop dan naar de kas,

van de landheer gepacht,
die hen meer voedsel bracht,

om fruit te verzorgen
of mest te bezorgen.

Dan was het bange wicht
uit haar ouders hun zicht.

Op een zomeravond
sloop Emmy schichtig rond

richting de glazen kas
toen jongens uit haar klas

plots achter haar stonden
en haar onomwonden

een kus wilden geven.
Zij ontweek hun streven.

Eentje greep Emmy vast
toen een andere gast

een zoen op haar wang gaf.
“Blijf liever van me af!”.

Doordat men echter zag
een giechelende lach

deed de zwijgende rest
niets met Emmy’s protest.

Ze kon zich verzoenen
met de verre zoenen.

Nog buitensporig was
een andere Judas.

Die bond haar middel vast
aan een vitrinekast

om kruiden te kweken.
Het glas niet te breken

van haar onderkomen
en toch los te komen

was geen enkel probleem,
doch verpestte de game.

De vastgebonden griet
verroerde zich dus niet.

Met haar vastgebonden
werd men opgewonden.

Zo werden onderwijl
de kwajongens bloedgeil.

Te jong om te weten
of passend vergeten

dat deze daad niet hoort.
De lust die de deugd smoort.

Te jong om te denken
aan wat ze kon schenken

aan lijfelijk genot.
De zoete honingpot

die onberoerd schuil ging
onder Emmy’s kleding.

De boys in het tuinhuis
streelden over hun kruis.

Kort waarna hand na hand
in een werkbroek beland.

Al eens gemasturbeerd
hadden ze zo geleerd

om vruchten te plukken
door zich af te rukken.

Zo gaven de ventjes
haar veel complimentjes.

“Ik zou je graag ruiken!”
“Ik wil op je duiken!”

“Ik wil alles wel zien!”
“Je bent een hete trien!”

Naïef keek ze verbaasd
hoe ze steeds meer gehaasd

hun verborgen trekhand
in krampachtige stand

in hun gulp verschoven.
Eén piemel kwam boven

een broekriem uitsteken
om daar uit te breken

door erlangs te slijpen.
De geschudde pijpen

pompten en rommelden.
De heupen schommelden.

De gespotloodventer
vond zijn drift urgenter,

getuige zijn snuiven,
dan het omhoog schuiven

van zijn verzakte broek.
Zijn decorum bleek zoek.

“Een hond die blaft, bijt niet”,
bleek een fout welkomstlied

van hun valse streken.
Grote ogen keken

naar de roze vulkaan
die wit magma liet gaan

en een hand besmeurde.
Wat tevens gebeurde
bij alle lulletjes
van de elf knulletjes;

maar onder hun gulpen.
De handen als hulpen

voor een kort hoogtepunt,
werd een zalfje gegund.

De snuiver niesde plots
en wilde zijn neuskots

blijkbaar onbevangen
in zijn hand opvangen.

Doch in dezelfde hand
waar zijn zaad was beland.

Zijn snot werd geborgen.
maar hij kreeg kopzorgen

door zijn snelle actie.
Een obscene tronie.

In ieders verbeelding
landde de ontlading

in het rode trutje.
Geen weet van haar kutje,

want dat kenden ze niet.
Er gold hetzelfde lied

voor de groene Emmy.
In haar lustfantasie

speelde zij hun handen.
Haar schaamlipjes brandden.

De blik van de schoffies,
het schuim uit hun lolly’s

en de lossende kreun
gaven allemaal steun

aan een nieuw roeseffect
dat vocht uit haar doos lekt.

Doch haar eroticum
kwam niettemin stiekem

voort uit ieders schaamte.
Dat was het geraamte

voor een latere schuld
die ze met veel geduld

geraffineerd inde,
wanneer het haar zinde.

Het voorkwam pesterij.
Emmy kneedde de klei

wanneer deze week was
bij jongens uit haar klas.

Plots stormde paps binnen.
Geheel buiten zinnen

schopte de wildebras
de jongens uit de kas

en bevrijdde Emmy
uit haar dwangpositie.

Het bleef niet bij die keer.
De dorpsjeugd vond steeds weer

voor de dociele meid
een trek gelegenheid.

Haar te overvallen
en haar uit te stallen
iedere keer plotsklaps.
Zij durfde niet haar paps

hiervan in te lichten.
Men zou haar betichten

dat zij hen heeft verleid.
Zo ging dat in die tijd.

“Word jij vaker gepest?”,
was een vraag die slechts test

op haar loyaliteit.
De eerbiedige meid

kon niet op de schouders
van haar tere ouders

haar pijnen neerleggen.
Ze kon dus niets zeggen.

Waarom het verbieden
dat de jongelieden

zich soms bevredigen
om zich te ledigen?

Ze kreeg van de abdis,
direct na de dagmis,

een uitleg waarom God
dit zag als een verbod.

Ze vroeg, nog niet bekeerd,
hoe ze hadden geleerd

zichzelf te verwennen?
Leerden ze dat kennen

van andere zielen?
“Deze acties vielen

onder aangeboren”,
kreeg Emmy te horen.

“Al die jonge knapen
zijn door God geschapen

en hiermee toegerust?",
vroeg Emmy doelbewust.

Toen viel de abdis stil.
Aftrekken was Gods wil!

Ze was haar onschuld kwijt.
Verloor haar waardigheid.

Angst voor verlies werd haat.
Emmy haatte de smaad

en haar machteloosheid.
“Maar als ik dit bestrijd,

het met weerzin verwerp
dan is mijn mes vlijmscherp

en controleer ik hen.
Zodat ik hoger ben!”

Werd ze superieur
of bleef toch de teneur

van een begeerd meisje
op haar beschaamd reisje?

Kreeg als Georges Bataille
Emmy represaille

door de expliciteit
van de moraliteit,

die hier werd geschonden?
Werd inzicht gevonden?
